\subsection{Preparing equipment and selling its first samples}
\label{sec:ppe-pre-use-cost}

Preparing a machine can produce two things of value: equipment capable of
operating as intended, and goods that can already be sold. Their costs and
receipts answer different accounting questions. Necessary delivery,
installation and functional testing contribute to the machine's cost under
AASB 116.16--17. Paragraph 20A requires the proceeds from selling items
produced during preparation, and those items' costs, to enter profit under
the applicable standards. It measures the goods' costs using AASB 102
\citep{aasb116dated,aasb102compiled2021}. The samples' margin therefore does
not reduce the cost of the machine.

Consider a synthetic Australian manufacturer in calendar 2024, starting with
cash and contributed capital of 300. It pays 100 for equipment, 5 for its
initial delivery, 12 for necessary installation and 3 for an engineering
service that tests and calibrates its physical performance. Probable benefits
and reliable cost measurement are admitted. These four amounts give PPE cost
of 120. The engineering service uses different resources from those recorded
as sample production below. Each receipt or resource has an identity and an
assigned purpose; sharing a resource would require an evidenced allocation,
not two recognitions of its full cost.

The sole sample batch uses normal materials costing 4 and directly attributable
contracted conversion costing 2. The example has no other production
overheads. Its goods cost is thus 6 under AASB 102.10--15. In a real factory,
paragraph 13 requires fixed production overheads to be allocated using normal
capacity; low output does not justify loading every unused resource onto a
few units. Unknown overheads would leave the goods' cost unresolved here.
Separate abnormal production waste of 1 is expensed under paragraph 16(a).
Staff training of 3 and advertising of 2 are also period expenses, excluded
from PPE by AASB 116.19. All these amounts are paid in cash.

On 30 June, the manufacturer sells the entire batch for cash of 8. These are
ordinary manufactured goods, outside the inventory standard's agricultural,
financial-instrument and broker-trader exceptions. Net realisable value is 8,
with no further completion or selling costs, so cost 6 needs no write-down.
The approved, enforceable customer contract identifies the goods and price,
has commercial substance and probable collection. Delivery, acceptance, title
and control transfer together on that date. There are no remaining promises,
returns or price concessions in this case. These facts admit the revenue
treatment under AASB 15.9 and 31--38; cash receipt alone would not establish it
\citep{aasb15dated}.

The entries preserve both assets. Acquisition and preparation debit PPE 120
and credit cash 120. Sample production debits inventory 6 and credits cash 6;
the excluded costs debit period expense 6 and credit cash 6. At sale, debit
cash 8 and credit revenue 8, then debit cost of sales 6 and credit inventory 6.
AASB 102.34 links the goods' expense to recognition of the related revenue.
The batch's cost history remains available even though its closing stock is
zero. Before depreciation, cash is 176, PPE is 120 and the loss is 4.

The machine becomes capable of its intended operation on 1 July but remains
idle through 31 December. No further production occurs in this example.
Management's admitted estimate is a five-year useful life, zero residual
value and uniform consumption of benefits, including obsolescence while idle.
The straight-line method uses equal monthly charges. Paragraph 55 starts
depreciation when the asset is available for use and does not suspend this
charge during idleness. It also allows a usage-based method to have zero
depreciation when production is zero; that is a different admitted consumption
pattern. Here the six available months give
\[
  \text{depreciation}=\frac{120-0}{60}\times6=12,
  \qquad \text{closing PPE}=120-12=108.
\]
Debit depreciation expense 12 and credit accumulated depreciation 12. Profit
is \(8-6-6-12=-16\), so equity closes at \(300-16=284\).
Cash 176 plus net PPE 108 gives the same total. The initial capital is already
in the opening balance; it is not a financing receipt in this year.

Cash-flow presentation follows the purchases and sale. PPE payments of 120
are investing outflows under AASB 107.16(a). The sample receipt of 8, goods
payments of 6 and period-cost payments of 6 give operating cash flow of
\(8-6-6=-4\); profit of \(-16\) plus noncash depreciation of 12 gives the
same result under the indirect method. Paragraphs 14 and 18--20 support this
classification and reconciliation \citep{aasb107dated}. Closing cash is
\(300-120-4=176\).

Several wrong treatments would still balance. Netting the samples' cost and
proceeds into PPE gives \(120+6-8=118\), depreciation 11.8, net PPE 106.2
and a loss of 17.8; cash remains 176. It also shifts the sample margin from
operating cash flow into investing cash flow. Waiting for productive use
would omit depreciation 12 and overstate both PPE and equity by 12.
Capitalising the excluded costs would instead produce cost 126 and closing
PPE 113.4. A balanced trial balance cannot distinguish these errors from the
required treatment. The method needs the cost purpose, goods population,
sale conditions, consumption pattern and availability date.

The four retained compilations have an overlapping application period that
includes this calendar 2024 example: AASB 102 compilation 3, AASB 116
compilation 6, AASB 15 compilation 8 and AASB 107 compilation 5. The example
assumes no early adoption and isolates the pretax cost and sale mechanism;
taxes, financing, grants, restoration obligations and later impairment require
their own facts. With initial cost established, replacing part of the machine
must next remove the cost history of the part that has gone.
